\section{项目背景与目标}

Grover搜索算法是量子计算中最具代表性的算法之一。对于含有\(N=2^n\)个候选状态、且只有一个目标解的无结构搜索问题，经典穷举搜索平均需要\(O(N)\)次查询，而Grover算法只需要约\(O(\sqrt{N})\)次Oracle查询即可显著提高目标态被测得的概率。

从应用角度看，Grover算法并不只对应“在数据库里找一条记录”这一种直观场景。更一般地，只要一个问题可以被表述为“在大量候选解中寻找满足判定条件的解”，就可以把判定条件封装为Oracle，再利用Grover式振幅放大提高正确候选被测量到的概率。典型应用包括：无结构数据库或索引检索、密码学中的密钥搜索与哈希原像搜索、组合优化中的可行解搜索、约束满足问题中的候选筛选，以及作为其他量子算法子程序的幅度放大模块。虽然真实应用还需要考虑Oracle构造成本、硬件噪声和量子资源开销，但Grover算法提供的平方级加速仍然是理解量子搜索优势的重要入口。

不过，对于初学者而言，Grover算法的难点往往不在最终测量结果，而在于理解“振幅放大”到底如何发生。Oracle本身并不会直接增加目标态概率，它只是把目标态的概率幅乘以\(-1\)，也就是完成相位标记。真正把相位标记转化为概率提升的是Diffusion算子，即关于平均概率幅的反演。

因此，本项目的目标设定为：
\begin{enumerate}
  \item 构造一个支持2到3个量子比特的Grover搜索玩具模型；
  \item 支持任意计算基目标态，例如\(|10\rangle\)、\(|101\rangle\)；
  \item 在Hadamard初始化、每次Oracle操作、每次Diffusion操作后保存中间statevector；
  \item 用2D柱状图和3D步进图展示各基态概率幅和概率的变化；
  \item 通过交互式滑块、动画和电路视图帮助观察“Oracle翻转”和“均值倒置”的作用差异。
\end{enumerate}
